\subsection{Multicast: One Request, a Structured Result}
\label{sec:multicast}

An ordinary send has one receiver. The original \filepath{lib/multicast.abclr2}
adds a library-defined operation: send the same message to every object in
a target tree, and optionally reconstruct a tree of replies. This is not
network multicast, a transaction, or the unsupported brace-delimited send
syntax. Its targets are ordinary ABCL/R2 objects in Lisp cons structures.
Where guards changed the receiving object's script selection, multicast
changes how the sender's evaluator interprets a target. Keep the future
reply-stream model and the evaluator context from earlier sections in mind.

\subsubsection{Replay the Original Library Example}

The library ends with a commented example, rather than a separate file in
\filepath{examples/}. Its three objects are:

\wholelisting[firstline=173,lastline=183]{lib/multicast.abclr2}

In a fresh session at the repository root, load the library, wrap the
top-level evaluator, enter those definitions, and replay the original send:

\begin{codeblock}
\begin{lstlisting}
(inc-lib "lib/multicast" :use-source t)
[*eval* := [multicast-eval-gen <== [:new *eval*]]]
; Enter the original twice, square, and identity definitions above.
[[twice square identity] <== 12]
; => (24 144 12)
\end{lstlisting}
\end{codeblock}

The bracket expression \code{[twice square identity]} constructs a list
of three object values. Each receives 12: \code{twice} returns 24,
\code{square} returns 144, and \code{identity} returns 12. The result list
uses the targets' positions, not the order in which replies arrive.
The same top-level installation is documented in the user's guide~\cite{guide}.
Loading the library alone only defines its machinery; it does not enable
list-valued targets in every evaluator.

\subsubsection{Follow the Shape, Not Just the Values}

Try a target tree containing both objects and literal data:

\begin{codeblock}
\begin{lstlisting}
[[twice [square :label] identity nil] <== 12]
; => (24 (144 :LABEL) 12 NIL)
[[twice twice] <== 3]                  ; => (6 6)
\end{lstlisting}
\end{codeblock}

Only object leaves receive the message. The keyword \code{:label} and
\code{nil} are retained as data; they are not receivers. Repeating an
object creates two target occurrences and two sends, not a deduplicated
recipient set. Each occurrence has its own result position.

The host Lisp function \code{pickup-target-objects} recursively walks
\code{car} and \code{cdr}. It constructs a fresh cons template with
\code{nil} in every object position, and records each object with a path
into that template. For the nested example:

\begin{longtable}{L{0.21\textwidth}L{0.37\textwidth}L{0.32\textwidth}}
\textbf{Receiver} & \textbf{Path from the template root} & \textbf{Stored reply}\\\hline
\code{twice} & \code{(:car)} & 24\\
\code{square} & \code{(:cdr :car :car)} & 144\\
\code{identity} & \code{(:cdr :cdr :car)} & 12\\
\end{longtable}

The template is \code{(nil (nil :label) nil nil)}. A path is a sequence
of operations: \code{(:cdr :car :car)} means take the outer tail, its
first nested list, and that list's first cell. The function builds its
target records with \code{push}, so their order is reversed relative to
the traversal. That does not change the reconstructed result.
\code{store-returned-value} follows a path and updates the final
\code{car} or \code{cdr} with \code{setf}.

This is a cons-tree walk, not arbitrary container traversal. A dotted
pair of object targets also works: \code{(cons twice square)} with message
3 yields \code{(6 . 9)}. Arrays are not recursively decomposed.
Use finite, acyclic target trees. A receiver that sends \code{!nil}
has replied; a \code{nil} result is not the same as a missing reply.

\subsubsection{Separate Submission from the Aggregate Reply}

The evaluator wrapper intercepts two ordinary send representations:

\wholelisting[firstline=134,lastline=143]{lib/multicast.abclr2}

For a nonempty cons target, a now-type request sends the evaluator's
continuation \code{C} as the aggregate reply destination. A past-type
request instead supplies its explicit application destination \code{Reply},
and preserves \code{C} separately as the submission continuation.
Single-object targets and other evaluator requests are forwarded to
\code{original-eval}. Future-type sends and extended send representations
are not intercepted by these two patterns; do not assume that every
extension composes automatically with multicast.

Read the \code{multicaster} handler in this order:

\begin{enumerate}[nosep]
  \item Decompose the target tree into template and object/path records.
  \item Send \code{nil} to its own continuation \code{C}, if present.
        This is submission progress, not a collected application result.
  \item Without an aggregate destination, send the payload to each object
        and collect nothing.
  \item Otherwise create one collector, initialize it with the template,
        target count, and aggregate destination, then send to each object
        with a position-specific reply receiver.
\end{enumerate}

\code{partial-value-receiver-gen} creates a \code{cont} object for each
position. An application reply becomes \code{[:result value position]}
to the collector. Notice another forwarding effect: the target's sender
binding is \code{multicaster}, not the original application client.
Sender-sensitive scripts must take this into account.

\subsubsection{Put a Multicast Client in Its Own Group}

The supplemental \filepath{tutorial/examples/10-multicast.abclr2}
keeps this policy local to a client group instead of replacing the
interactive evaluator:

\wholelisting[firstline=4,lastline=17]{tutorial/examples/10-multicast.abclr2}

The targets may remain ordinary objects: it is the \emph{sender's}
evaluator that interprets a list-valued target. The internal
\code{r2sys::create-primary-eval} helper is the same implementation-specific
helper used earlier. In a fresh session, enter the original three object
definitions, then load the supplemental file:

\begin{codeblock}
\begin{lstlisting}
(inc "tutorial/examples/10-multicast" :use-source t)
[multicast-client <== [:now [twice square] 3]]
; => (6 9)
[multicast-client
 <== [:past [twice square] 3 multicast-inbox]]
; => :SUBMITTED
[multicast-inbox <== :status]           ; => (1 (6 9))
\end{lstlisting}
\end{codeblock}

The inbox stores an arrival count and the last aggregate. Inspect its
short definition in the supplemental file. \code{:submitted} is the
client handler's reply to the interactive caller; \code{(6 9)} travels
separately to the inbox. The submission reply does not acknowledge that
every recipient has finished. The checked session observes the aggregate
after the finite sends have run; it does not establish an arrival-time
guarantee for other schedules.

For an existing group, the original \code{multicast-installer} accepts
\code{[:install group-manager]} as a past-type request. The manager must
support \code{:replace-evaluator}. Install before creating clients that
will use multicast: replacing a group's evaluator does not automatically
rewrite every existing object's captured evaluator. The guide's
\code{*eval*} assignment is the explicit top-level alternative.

\subsubsection{Know What the Collector Can and Cannot Promise}

The essential collector handlers are:

\wholelisting[firstline=72,lastline=82]{lib/multicast.abclr2}

Initialization stores the template, remaining-result count, and reply
destination. Here \code{C} names the numeric count, not the evaluator
continuation used earlier; \code{R} is the aggregate destination.
Each \code{:result} stores a value and decrements the count.
Only the transition to zero sends the aggregate. This assumes
\emph{exactly one reply per target occurrence}.

A nonempty target list with no object leaves initializes a count of zero,
but initialization sends no aggregate. A silent recipient leaves the
count above zero. Neither case has a library timeout or failure result;
a now-type caller cannot obtain the missing aggregate. An empty list is
\code{nil}, not a cons, so it does not enter the multicast branch at all.
Use a nonempty set of object recipients for collected requests, and make
each recipient's one-reply contract explicit.

The collector does not remember which positions have already replied.
A controlled probe initializes two positions, supplies two results for
the first position, and observes a premature \code{(9 nil)} aggregate.
A later result for the second position mutates that same delivered list
to \code{(9 10)} without a second aggregate message. These are observed
protocol limits, not a supported duplicate-delivery or retry mechanism.
Do not use a multi-reply stream producer as a one-reply multicast recipient.
For fire-and-forget sends, replies are not collected and there is no
aggregate completion acknowledgment.

\paragraph{Evidence and exercise.}
Run \code{make -C tutorial multicast-check}. Twenty-four checks replay
the original commented session and test nested/dotted targets, duplicate
occurrences, \code{nil} replies, sender identity, group setup, explicit
reply routing, fire-and-forget effects, and the collector boundaries above.
The runner extracts original lines 171--185 and copies the library
byte-for-byte; inputs and transcripts stay in \filepath{build/multicast-tests/}.
No original source or runtime is patched.

Why does \code{:label} survive in the result? It occupies a literal
template position rather than an object position. Why are two instances
of \code{twice} not merged? Positions, not distinct identities, define
the requests. As an exercise, add a one-reply negation object to a nested
target tree, predict the template and paths, then compare the collected
result. Keep the original library unchanged.
